\chapter{RAG -- Retrieval-Augmented Generation in der Praxis}
\label{chap:rag_vector_dbs}

% ============================================================
% LERNZIELE
% ============================================================
\begin{lernziele}
Nach diesem Kapitel kannst du:
\begin{itemize}
    \item Die vier Phasen einer RAG-Pipeline (Indexing, Retrieval, Augmentation, Generation) im Detail erklären
    \item Chunking-Strategien für verschiedene Dokumenttypen auswählen und implementieren
    \item Embedding-Modelle und Vector Databases anhand konkreter Anforderungen bewerten
    \item Eine produktionsreife RAG-Pipeline mit LangChain / LlamaIndex aufbauen
    \item Fortgeschrittene RAG-Techniken (HyDE, Multi-Query, RAPTOR) gezielt einsetzen
    \item RAG-Qualität mit Metriken (Hit Rate, MRR, Faithfulness) evaluieren
\end{itemize}
\end{lernziele}

% ============================================================
% MOTIVATION
% ============================================================
\section{Motivation}

Pre-trained Models haben zwei fundamentale Limitationen für Business-Anwendungen:

\begin{enumerate}[leftmargin=*]
    \item \textbf{Kein Firmenwissen}: Modelle kennen keine internen Dokumente, Produktdatenbanken oder Kundendaten deines Unternehmens.
    \item \textbf{Veraltete Information}: Der Knowledge Cutoff ist ein reales Problem -- Modelle wissen nichts über aktuelle Events, Produkte oder Preisänderungen.
\end{enumerate}

Fine-Tuning könnte beide Probleme lösen, ist aber teuer, aufwendig und unflexibel: Ein fine-getuntes Modell kann nur das Wissen, das zum Zeitpunkt des Trainings existierte. Neue Produkte oder aktualisierte Richtlinien bedeuten ein erneutes Training.

\textbf{RAG (Retrieval-Augmented Generation)} löst beide Probleme grundlegend anders: Statt Wissen in den Modellgewichten zu speichern, wird es bei jeder Anfrage aus einer externen Wissensdatenbank abgerufen und als Kontext in den Prompt eingefügt. Das Ergebnis sind aktuelle, faktenbasierte Antworten -- ohne Fine-Tuning, ohne Modell-Updates, mit vollständiger Kontrolle über die Wissensbasis.

RAG ist 2026 das dominierende Pattern für LLM-Produktionssysteme: über 70\,\% aller produktiven LLM-Deployments nutzen RAG.

% ============================================================
% GRUNDLAGEN
% ============================================================
\section{Grundlagen -- RAG-Konzepte}

\subsection{Das Problem, das RAG löst}

\begin{verbatim}
   Ohne RAG:
   [User-Frage] --> [LLM] --> [Antwort (nur aus Training)]
   
   Mit RAG:
   [User-Frage] --> [Retriever] --> [Vector DB]
                       |                |
                       v                v
   [Retrieved Docs] + [User-Frage] --> [LLM] --> [Faktenbasierte Antwort]
\end{verbatim}

\subsection{Die vier Phasen von RAG}

\begin{description}[leftmargin=*,style=nextline]
    \item[Indexing (Indexierung):] Dokumente werden in kleine Chunks zerschnitten und jedes Chunk wird in einen Vektor (Embedding) umgewandelt. Dieser Vektor landet in einer Vector Database.
    \item[Retrieval (Abruf):] Bei einer User-Anfrage wird auch diese Anfrage in einen Vektor umgewandelt und im Index nach den ähnlichsten Chunks gesucht -- per Cosine Similarity oder Maximal Marginal Relevance (MMR).
    \item[Augmentation (Anreicherung):] Die gefundenen Chunks werden zum Prompt hinzugefügt -- als Kontext, auf den das Modell zugreifen kann.
    \item[Generation (Generierung):] Das LLM liest den Prompt inklusive der Context-Chunks und generiert eine Antwort, die auf den retrieveden Fakten basiert und diese zitiert.
\end{description}

\begin{verbatim}
   Indexing-Phase:
   [PDF] -> [Chunker] -> [Chunks] -> [Embedding] -> [Vector DB]
                                         |
                                   text-embedding-3-small
   
   Runtime-Phase:
   [Query] -> [Embedding] -> [Vector DB (ANN Search)] -> [Top-K Chunks]
                                                              |
                                                              v
   [System-Prompt] + [Chunks] + [Query] -> [LLM] -> [Antwort]
\end{verbatim}

\subsection{Naive RAG vs. Advanced RAG}

\begin{table}[H]
\centering
\begin{tabularx}{\linewidth}{lX}
\toprule
\textbf{Typ} & \textbf{Ansatz} \\
\midrule
Naive RAG & Einfaches Chunking + Cosine Similarity + Top-K Retrieval \\
Advanced RAG & Multi-Query, HyDE, RAPTOR, Corrective RAG, Self-RAG \\
\midrule
\end{tabularx}
\caption{Naive vs. Advanced RAG}
\label{tab:rag-types}
\end{table}

% ============================================================
% ARCHITEKTUR
% ============================================================
\section{Architektur -- RAG-Systeme in Produktion}

\subsection{Komponenten einer RAG-Produktion}

\begin{verbatim}
   +-------------------+     +-------------------+     +-------------------+
   |  Data Ingestion   | --> |  Embedding + Index | --> |  Vector Database  |
   |  (PDFs, DB, Web)  |     |  (asynchron)       |     |  (pgvector/Pinecone)|
   +-------------------+     +-------------------+     +-------------------+
                                                              |
                                                              v
   +-------------------+     +-------------------+     +-------------------+
   |  User Query       | --> |  Embedder + Retrieval| -->|  LLM + Response   |
   |  (REST API)       |     |  (ANN Search)       |     |  (GPT-4o/Claude)  |
   +-------------------+     +-------------------+     +-------------------+
\end{verbatim}

\subsection{Indexing-Pipeline (Offline)}

Die Indexing-Pipeline läuft asynchron -- bei Initialisierung, bei Datenänderung oder auf einem Cron-Job:

\begin{enumerate}[leftmargin=*]
    \item \textbf{Dokumenten-Loader}: PDF via PyMuPDF, HTML via Crawler, Markdown aus Git-Repo, SQL aus Datenbank
    \item \textbf{Pre-Processing}: Metadaten extrahieren (Titel, Datum, Autor), HTML-Tags entfernen, Tabellen in Text umwandeln
    \item \textbf{Chunking}: Recursive Splitting nach Absätzen, Sätzen, Wörtern
    \item \textbf{Embedding}: Jeder Chunk wird embeddet (OpenAI / Cohere / Voyage)
    \item \textbf{Insert}: Embedding + Metadaten + Originaltext in Vector DB schreiben
\end{enumerate}

\subsection{Runtime-Pipeline (Online)}

Die Runtime-Pipeline läuft bei jeder User-Anfrage:

\begin{enumerate}[leftmargin=*]
    \item \textbf{Query Embedding}: Die User-Frage wird mit dem \textit{gleichen} Embedding-Modell in einen Vektor umgewandelt
    \item \textbf{ANN Search}: Approximate Nearest Neighbor -- die "nächsten" Vektoren im Index finden
    \item \textbf{Re-Ranking}: Optional -- die Top-K-Ergebnisse durch ein Cross-Encoder-Modell neu sortieren
    \item \textbf{Prompt Construction}: System-Prompt + Retrieved Chunks + User-Frage zusammenbauen
    \item \textbf{Generation}: LLM generiert Antwort basierend auf dem Kontext
\end{enumerate}

% ============================================================
% THEORIE
% ============================================================
\section{Theorie -- Warum RAG funktioniert (und wann nicht)}

\subsection{Embedding als semantischer Raum}

Ein Embedding-Modell projiziert Text in einen hochdimensionalen Vektorraum. Die entscheidende Eigenschaft: \textbf{Semantisch ähnliche Texte liegen in diesem Raum nahe beieinander}.

Die Distanz zwischen zwei Texten wird meist via \textbf{Cosine Similarity} gemessen:

\[
\operatorname{cosine\_similarity}(A, B) = \frac{A \cdot B}{\|A\| \cdot \|B\|}
\]

Je näher der Wert an 1, desto ähnlicher sind die Texte. Ein Wert von 0.8+ bedeutet hohe semantische Ähnlichkeit, 0.3 bedeutet kaum Zusammenhang.

\subsection{Das Vokabular-Problem}

Embeddings sind gut für semantische Ähnlichkeit, aber \textbf{schlecht für exakte Wortübereinstimmungen}. Ein Chunk, der "Kühlschrank" enthält, wird nicht gefunden, wenn die Query "Kühlgerät" ist -- auch wenn es semantisch nahe liegt. Zwei Chunks mit "Kühlschrank" und "Waschmaschine" können ähnliche Embeddings haben, wenn sie im selben Dokumentkontext stehen.

\textbf{Lösung}: Hybride Suche (BM25 + Embedding) kombiniert die Stärken beider Ansätze.

\subsection{Das Relevanz-Problem}

Nicht jedes ähnliche Chunk ist auch relevant für die Beantwortung einer Frage. Cosine Similarity misst Text-Ähnlichkeit, nicht Question-Answering-Relevanz. Ein Chunk über "Rückgabepolitik" und eine Frage zu "Rückgabepolitik" haben hohe Similarity -- aber ein Chunk über "Versandkosten" kann für die Frage "Wie hoch sind die Versandkosten?" trotzdem irrelevant sein, wenn das Chunk nur "Versandkosten sind im Preis enthalten" sagt.

\textbf{Lösung}: Cross-Encoder Re-Ranker verbessern die Retrieval-Qualität signifikant (siehe Abschnitt~\ref{sec:reranker}).

\subsection{Metriken für RAG-Qualität}

\begin{table}[H]
\centering
\begin{tabularx}{\linewidth}{lX}
\toprule
\textbf{Metrik} & \textbf{Beschreibung} \\
\midrule
Hit Rate @ K & Wie oft ist die relevante Information in den Top-K? \\
MRR (Mean Reciprocal Rank) & Wie weit oben im Ranking ist die erste relevante Information? \\
NDCG @ K & Berücksichtigt die Position aller relevanten Ergebnisse \\
Faithfulness & Wie viel der LLM-Antwort ist durch retrievede Chunks belegbar? \\
Answer Relevance & Wie relevant ist die Antwort für die Frage? \\
\midrule
\end{tabularx}
\caption{RAG-Qualitätsmetriken}
\label{tab:rag-metrics}
\end{table}

% ===========================================================%
% CHUNKING (existiert, erweitert)
% ===========================================================%
\section{Chunking-Strategien}

Die Chunking-Strategie ist der \textbf{wichtigste Hebel für RAG-Qualität}. Schlecht gechunkte Dokumente führen unweigerlich zu schlechten Retrieval-Ergebnissen.

\subsection{Strategien im Detail}

\begin{table}[H]
\centering
\begin{tabularx}{\linewidth}{lXX}
\toprule
\textbf{Strategie} & \textbf{Vorteil} & \textbf{Nachteil} \\
\midrule
Fixed-Size (256-512 Tokens) & Einfach, performant & Kontextbrüche an Grenzen \\
Semantic Chunking & Thema bleibt intakt & Langsamer, komplexer \\
Recursive Chunking & Struktur-Integrität & Overhead bei kurzen Texten \\
Agentic Chunking (LLM-gesteuert) & Optimal pro Dokument & Teuer (LLM-Calls) \\
Late Chunking & Bessere Kontexterhaltung & Experimentell \\
\midrule
\end{tabularx}
\caption{Chunking-Strategien im Vergleich}
\label{tab:chunking}
\end{table}

\subsection{Chunking implementieren}

\begin{lstlisting}[language=Python, caption=Verschiedene Chunking-Strategien]
from langchain.text_splitter import (
    RecursiveCharacterTextSplitter,
    TokenTextSplitter,
    MarkdownHeaderTextSplitter,
)

class ChunkingStrategy:
    @staticmethod
    def recursive(docs: list[str], chunk_size: int = 512,
                  overlap: int = 64) -> list[str]:
        splitter = RecursiveCharacterTextSplitter(
            chunk_size=chunk_size,
            chunk_overlap=overlap,
            separators=["\n\n", "\n", ".", " ", ""],
        )
        return splitter.split_documents(docs)

    @staticmethod
    def markdown_headers(docs: list[str]) -> list[str]:
        splitter = MarkdownHeaderTextSplitter(
            headers_to_split_on=[
                ("#", "H1"),
                ("##", "H2"),
                ("###", "H3"),
            ]
        )
        return splitter.split_text(docs[0])

    @staticmethod
    def semantic(chunks: list[str],
                 model: str = "gpt-4o-mini") -> list[str]:
        result = []
        current = ""
        for chunk in chunks:
            if current and len(current) > 256:
                result.append(current)
                current = chunk
            else:
                current += "\n" + chunk
        if current:
            result.append(current)
        return result

    @staticmethod
    def agentic(doc: str) -> list[str]:
        response = openai_client.chat.completions.create(
            model="gpt-4o-mini",
            messages=[{
                "role": "system",
                "content": (
                    "Teile das Dokument in semantische Abschnitte. "
                    "Gib JSON: [{'content': '...', 'heading': '...'}]"
                )
            }, {
                "role": "user",
                "content": doc[:8000],
            }],
            response_format={"type": "json_object"},
        )
        return json.loads(response.choices[0].message.content)
\end{lstlisting}

\subsection{Empfehlungen}

\begin{itemize}[leftmargin=*]
    \item \textbf{Starte mit Recursive Chunking} (512 Tokens, 64 Overlap) -- es funktioniert für 80\,\% der Fälle
    \item \textbf{Füge Metadaten hinzu}: Document-ID, Heading-Hierarchie, Seitenzahl. Das verbessert das Retrieval.
    \item \textbf{Overlap ist Pflicht}: 10--20\,\% Overlap verhindert Informationsverlust an Chunk-Grenzen
    \item \textbf{Document-scoped retrieval}: Nach dem Retrieval die "Nachbar-Chunks" desselben Dokuments mitladen, die an der Grenze liegen
\end{itemize}

% ===========================================================%
% EMBEDDING-MODELLE
% ===========================================================%
\section{Embedding-Modelle im Detail}

\subsection{Embedding-Modelle im Vergleich}

\begin{table}[H]
\centering
\begin{tabularx}{\linewidth}{lXrr}
\toprule
\textbf{Modell} & \textbf{Besonders gut für} & \textbf{Dimension} & \textbf{Kosten/1K Tokens} \\
\midrule
text-embedding-3-large & Allgemein, hohe Genauigkeit & 3072 & \$0,00013 \\
text-embedding-3-small & Allgemein, günstig & 1536 & \$0,00002 \\
Cohere embed-multilingual-v3 & Multilingual (Deutsch, 100+ Sprachen) & 1024 & Kostenlos (Trial) \\
Voyage-2-large & Finance, Legal, Code & 1024 & \$0,00006 \\
BAAI/bge-m3 (Open Source) & Selbst gehostet, multilingual & 1024 & Kostenlos \\
intfloat/multilingual-e5-large & Hohe Qualität, selbst gehostet & 1024 & Kostenlos \\
\midrule
\end{tabularx}
\caption{Embedding-Modelle 2026 im Vergleich}
\label{tab:embeddings}
\end{table}

\subsection{Praktische Embedding-Implementierung}

\begin{lstlisting}[language=Python, caption=Embedding-Service mit Multi-Provider-Unterstützung]
from openai import OpenAI
import numpy as np

class EmbeddingService:
    def __init__(self, provider: str = "openai"):
        self.provider = provider
        self.client = OpenAI()

    def embed(self, texts: list[str]) -> list[list[float]]:
        if self.provider == "openai":
            response = self.client.embeddings.create(
                model="text-embedding-3-small",
                input=texts,
            )
            return [r.embedding for r in response.data]

        elif self.provider == "cohere":
            import cohere
            co = cohere.Client()
            response = co.embed(
                texts=texts,
                model="embed-multilingual-v3.0",
                input_type="search_document",
            )
            return response.embeddings

        elif self.provider == "voyage":
            import voyageai
            vo = voyageai.Client()
            response = vo.embed(
                texts=texts,
                model="voyage-2",
                input_type="document",
            )
            return response.embeddings

    def similarity(self, a: list[float], b: list[float]) -> float:
        a_np = np.array(a)
        b_np = np.array(b)
        return float(np.dot(a_np, b_np) / (
            np.linalg.norm(a_np) * np.linalg.norm(b_np)
        ))
\end{lstlisting}

% ===========================================================%
% VECTOR DATABASES
% ===========================================================%
\section{Vector Databases im Vergleich}

\subsection{Die wichtigsten Optionen}

\begin{table}[H]
\centering
\begin{tabularx}{\linewidth}{lXXX}
\toprule
\textbf{Datenbank} & \textbf{Typ} & \textbf{Index-Typ} & \textbf{Skalierung} \\
\midrule
Pinecone & Managed Service & Hierarchical Navigable Small World (HNSW) & Auto-Skalierung \\
pgvector & PostgreSQL Extension & IVFFlat, HNSW & Horizontal (PostgreSQL-Cluster) \\
ChromaDB & Local / Embedded & HNSW & Single-Node \\
Weaviate & Self-hosted / Cloud & HNSW + BM25 Hybrid & Kubernetes \\
Qdrant & Self-hosted / Cloud & HNSW + Product Quantization & Sharding \\
Milvus & Distributed & IVF, HNSW, DiskANN & 100B+ Vektoren \\
\midrule
\end{tabularx}
\caption{Vector Databases 2026 im Vergleich}
\label{tab:vectordbs}
\end{table}

\subsection{Entscheidungsmatrix}

\begin{itemize}[leftmargin=*]
    \item \textbf{Startup / Prototyp}: ChromaDB (kein Setup, keine Kosten)
    \item \textbf{PostgreSQL-Shop}: pgvector (keine neue Infrastruktur)
    \item \textbf{Production / Scale}: Pinecone oder Qdrant (Managed, zuverlässig)
    \item \textbf{Enterprise / Multi-Cloud}: Weaviate (Hybrid-Suche, GraphQL)
    \item \textbf{Extrem-Scale (100M+ Vektoren)}: Milvus / Zilliz
\end{itemize}

\subsection{pgvector in der Praxis}

\begin{lstlisting}[language=Python, caption=RAG mit pgvector und SQLAlchemy]
import os
from sqlalchemy import create_engine, text
from pgvector.sqlalchemy import Vector

class PgVectorStore:
    def __init__(self, connection_string: str = None):
        self.engine = create_engine(
            connection_string or os.environ["DATABASE_URL"]
        )
        self._init_vector_extension()

    def _init_vector_extension(self):
        with self.engine.connect() as conn:
            conn.execute(text("CREATE EXTENSION IF NOT EXISTS vector"))
            conn.execute(text("""
                CREATE TABLE IF NOT EXISTS document_chunks (
                    id SERIAL PRIMARY KEY,
                    content TEXT,
                    metadata JSONB,
                    embedding vector(1536)
                )
            """))
            conn.execute(text("""
                CREATE INDEX IF NOT EXISTS idx_embedding
                ON document_chunks
                USING hnsw (embedding vector_cosine_ops)
                WITH (m = 16, ef_construction = 200);
            """))
            conn.commit()

    def insert(self, content: str, embedding: list[float],
               metadata: dict = None):
        with self.engine.connect() as conn:
            conn.execute(
                text("""
                    INSERT INTO document_chunks
                    (content, embedding, metadata)
                    VALUES (:content, :embedding, :metadata)
                """),
                {
                    "content": content,
                    "embedding": embedding,
                    "metadata": metadata or {},
                }
            )
            conn.commit()

    def search(self, query_embedding: list[float],
               top_k: int = 5) -> list[dict]:
        with self.engine.connect() as conn:
            result = conn.execute(
                text("""
                    SELECT content, metadata,
                           1 - (embedding <=> :query) AS similarity
                    FROM document_chunks
                    ORDER BY embedding <=> :query
                    LIMIT :top_k
                """),
                {"query": query_embedding, "top_k": top_k}
            )
            return [
                {
                    "content": row[0],
                    "metadata": row[1],
                    "similarity": row[2],
                }
                for row in result
            ]
\end{lstlisting}

% ===========================================================%
% FORTGESCHRITTENE RAG-TECHNIKEN
% ===========================================================%
\section{Fortgeschrittene RAG-Techniken}

\subsection{HyDE (Hypothetical Document Embeddings)}

Statt die Query direkt zu embedden, wird zuerst eine hypothetische Antwort generiert und \textit{diese} embeddet. Die Idee: Eine Antwort enthält mehr relevante Keywords und Kontext als eine kurze Frage.

\begin{lstlisting}[language=Python, caption=HyDE -- bessere Retrieval-Qualität]
def hyde_retrieval(query: str, retriever, llm) -> list[str]:
    """Generiert eine hypothetische Antwort und nutzt sie als Query."""

    # Schritt 1: Hypothetische Antwort generieren
    hyde_response = llm.invoke(
        f"Beantworte die folgende Frage knapp: {query}"
    )

    # Schritt 2: Die Antwort embedden (statt der Query)
    hyde_embedding = embedding_service.embed([hyde_response])[0]

    # Schritt 3: Mit dem Embedding suchen
    results = vector_db.search(hyde_embedding, top_k=5)
    return [r["content"] for r in results]
\end{lstlisting}

\subsection{Multi-Query Retrieval}

Statt einer Query werden mehrere Varianten generiert, um verschiedene Aspekte abzudecken:

\begin{lstlisting}[language=Python, caption=Multi-Query Retrieval]
def multi_query_retrieval(query: str, retriever, llm,
                          n_queries: int = 3) -> list[str]:
    queries = llm.invoke(
        f"Generiere {n_queries} verschiedene Formulierungen "
        f"dieser Frage, um alle relevanten Informationen zu finden:"
        f"\n{query}"
    )

    all_docs = []
    for q in queries:
        docs = retriever.get_relevant_documents(q)
        all_docs.extend(docs)

    # Deduplizieren
    seen = set()
    unique = []
    for d in all_docs:
        if d.page_content not in seen:
            seen.add(d.page_content)
            unique.append(d)
    return unique[:top_k]
\end{subtract}
\end{lstlisting}

\subsection{RAPTOR (Recursive Abstractive Processing)}

RAPTOR baut eine Hierarchie von Zusammenfassungen auf -- Textblöcke werden zusammengefasst, diese Zusammenfassungen wiederum. Das Retrieval sucht auf der passenden Hierarchie-Ebene:

\begin{verbatim}
   Level 2: [Gesamt-Zusammenfassung]
                /        \
   Level 1: [Summary A]  [Summary B]  ...
               |            |
   Level 0: [Chunk 1-5]  [Chunk 6-10]  ...
   
   Vorteil: Uebergreifende Fragen werden auf Level 1/2 beantwortet,
   Detail-Fragen auf Level 0.
\end{verbatim}

\subsection{Corrective RAG (CRAG)}

CRAG fügt einen \textbf{Qualitäts-Check} zwischen Retrieval und Generation ein:

\begin{enumerate}[leftmargin=*]
    \item Retriever liefert Top-K Chunks
    \item Ein Evaluator-Modell prüft: Sind die Chunks relevant? (Score 0--1)
    \item \textbf{Wenn relevant}: Chunks werden an LLM übergeben (normaler RAG-Flow)
    \item \textbf{Wenn nicht relevant}: Fallback -- LLM soll "Dokumenten nicht gefunden" antworten
    \item \textbf{Wenn unsicher}: Web-Suche als zusätzliche Quelle nutzen
\end{enumerate}

\subsection{Self-RAG}

Self-RAG lässt das LLM selbst entscheiden, ob und wann es retrievalt:

\begin{enumerate}[leftmargin=*]
    \item LLM erhält Query + Option zum Retrieval
    \item Wenn LLM Retrieval anfordert: Chunks werden eingespielt
    \item LLM markiert, welche Teile der Antwort durch Chunks belegt sind (\texttt{citation}-Marker)
    \item Ein zweites Modell prüft die Korrektheit der Zitationen
\end{enumerate}

% ===========================================================%
% PRAXIS: KOMPLETTE RAG-PIPELINE
% ===========================================================%
\section{Praxis -- Komplette RAG-Pipeline}

\begin{lstlisting}[language=Python, caption=Produktionsreife RAG-Pipeline]
import os
from dataclasses import dataclass, field
from typing import Optional

@dataclass
class RAGConfig:
    embedding_model: str = "text-embedding-3-small"
    llm_model: str = "gpt-4o-mini"
    chunk_size: int = 512
    chunk_overlap: int = 64
    top_k: int = 5
    rerank: bool = False
    hyde: bool = False

class RAGPipeline:
    def __init__(self, config: RAGConfig):
        self.config = config
        self.embedder = EmbeddingService(provider="openai")
        self.vector_store = self._init_vector_store()
        self.llm = openai_client

    def _init_vector_store(self):
        return Chroma(
            collection_name="rag_knowledge_base",
            embedding_function=OpenAIEmbeddings(
                model=self.config.embedding_model
            ),
            persist_directory="./chroma_db",
        )

    def add_documents(self, documents: list[str],
                      metadata: Optional[list[dict]] = None):
        chunks = ChunkingStrategy.recursive(
            documents,
            chunk_size=self.config.chunk_size,
            overlap=self.config.chunk_overlap,
        )
        metadatas = metadata or [{}] * len(chunks)
        self.vector_store.add_documents(chunks, metadatas=metadatas)

    def query(self, question: str) -> dict:
        if self.config.hyde:
            hyde_response = self.llm.invoke(
                f"Beantworte knapp: {question}"
            )
            query_text = hyde_response
        else:
            query_text = question

        docs = self.vector_store.similarity_search(
            query_text,
            k=self.config.top_k,
        )

        if self.config.rerank:
            docs = self._rerank(question, docs)

        context = "\n\n---\n\n".join(d.page_content for d in docs)
        response = self.llm.chat.completions.create(
            model=self.config.llm_model,
            messages=[{
                "role": "system",
                "content": (
                    "Du beantwortest Fragen basierend auf dem Kontext. "
                    "Wenn der Kontext nicht ausreicht: 'Nicht genug Daten.' "
                    "Zitiere Quellen aus dem Kontext."
                )
            }, {
                "role": "user",
                "content": f"Kontext:\n{context}\n\nFrage:\n{question}",
            }],
            temperature=0.1,
        )

        return {
            "answer": response.choices[0].message.content,
            "sources": [
                {"content": d.page_content, "score": d.metadata.get("score")}
                for d in docs
            ],
            "usage": response.usage,
        }
\end{lstlisting}

\subsection{Re-Ranking mit Cross-Encoders}
\label{sec:reranker}

\begin{lstlisting}[language=Python, caption=Re-Ranking verbessert Retrieval-Qualität]
class CrossEncoderReranker:
    def __init__(self, model_name: str = "cross-encoder/ms-marco-MiniLM-L-6-v2"):
        from sentence_transformers import CrossEncoder
        self.model = CrossEncoder(model_name)

    def rerank(self, query: str, documents: list[str],
               top_k: int = 3) -> list[tuple[str, float]]:
        pairs = [[query, doc] for doc in documents]
        scores = self.model.predict(pairs)
        scored = sorted(
            zip(documents, scores),
            key=lambda x: x[1],
            reverse=True,
        )
        return scored[:top_k]
\end{lstlisting}

% ===========================================================%
% BEST PRACTICES
% ===========================================================%
\section{Best Practices}

\begin{itemize}[leftmargin=*]
    \item \textbf{Hybride Suche verwenden}: Kombiniere Embedding-Similarity mit BM25-Keyword-Search (z.\,B. Weaviate, Elasticsearch + pgvector). Fängt Exakt-Matches besser als reine Embeddings.
    \item \textbf{Chunk-Größe evaluieren}: Starte mit 512 Tokens, messe Hit Rate @ 5 und erhöhe/reduziere systematisch. Keine allgemeingültige "beste" Größe.
    \item \textbf{Embedding-Modell zum Retrieval testen}: Ein multilinguales Modell (Cohere, bge-m3) kann für deutsche Dokumente besser sein als OpenAI.
    \item \textbf{Metadaten-Filter nutzen}: Filtere nach Datum, Kategorie oder Document-Type vor der Similarity-Suche. Das reduziert False Positives drastisch.
    \item \textbf{Re-Ranking ist Pflicht ab 10.000+ Chunks}: Der retriever findet gute Chunks, aber der Re-Ranker sortiert sie optimal für den LLM-Kontext.
    \item \textbf{Index regelmäßig neu bauen}: Bei Datenänderung muss der Index aktualisiert werden. Implementiere ein Update-Skript, das läuft, wenn sich die Quelle ändert.
    \item \textbf{Prompt-Caching für System-Prompts}: Der RAG-System-Prompt (Instructions + Context-Vorlage) ändert sich selten -- cache ihn.
\end{itemize}

% ===========================================================%
% ANTI-PATTERNS
% ===========================================================%
\section{Anti-Patterns}

\begin{itemize}[leftmargin=*]
    \item \textbf{Chunks zu groß}: 1024+ Tokens pro Chunk verlieren semantische Präzision. Die Antwort wird vage, weil das Chunk zu viele Themen enthält.
    \item \textbf{Kein Chunk-Overlap}: Ohne Overlap (10--20\,\%) gehen Informationen an Chunk-Grenzen verloren -- besonders bei Absatzübergängen.
    \item \textbf{Nur Cosine Similarity}: Reine semantische Suche versagt bei Produktnamen, IDs und spezifischen Begriffen. Nutze hybride Suche.
    \item \textbf{Embedding-Modell nie aktualisieren}: Embedding-Modelle werden besser. Bleibe auf dem aktuellen Stand und migriere bei Bedarf (mit Re-Indexing).
    \item \textbf{Keine Evaluierung}: Ohne Hit Rate, MRR oder Faithfulness weißt du nicht, ob dein RAG-System gut ist. Definiere Metriken, bevor du optimierst.
    \item \textbf{Alle Chunks gleich behandeln}: Ein 200-Zeilen-Vertrag braucht andere Chunking-Strategie als ein 2-Zeilen-FAQ. Verschiedene Dokumenttypen = verschiedene Strategien.
    \item \textbf{Keine Quellenangaben in der Antwort}: Der LLM sollte immer die Quelle zitieren (z.\,B. "Laut Dokument X auf Seite Y..."). Sonst sind die Antworten nicht überprüfbar.
\end{itemize}

% ===========================================================%
% PERFORMANCE
% ===========================================================%
\section{Performance und Optimierung}

\subsection{Latenz-Budget einer RAG-Query}

\begin{table}[H]
\centering
\begin{tabularx}{\linewidth}{lXr}
\toprule
\textbf{Schritt} & \textbf{Dauer} & \textbf{Anteil} \\
\midrule
Query Embedding & 50--200ms & 5--10\,\% \\
ANN Search (10K Vektoren) & 20--50ms & 2--5\,\% \\
ANN Search (1M Vektoren) & 100--500ms & 10--30\,\% \\
Prompt Construction & <1ms & <1\,\% \\
LLM Generation & 500--3000ms & 50--80\,\% \\
\midrule
\end{tabularx}
\caption{Latenz-Budget einer typischen RAG-Query}
\label{tab:rag-latency}
\end{table}

\subsection{Optimierungs-Strategien}

\begin{itemize}[leftmargin=*]
    \item \textbf{ANN statt brute-force}: Approximate Nearest Neighbor (HNSW-Index) statt exakter Suche -- 10x schneller bei <1\,\% Genauigkeitsverlust.
    \item \textbf{Vektordimension reduzieren}: text-embedding-3-large (3072) $\rightarrow$ text-embedding-3-small (1536) = 2x schnellere Suche bei minimalem Qualitätsverlust.
    \item \textbf{Quantisierung}: Product Quantization (PQ) reduziert Vektor-Speicher um 4--8x bei geringem Genauigkeitsverlust.
    \item \textbf{Pre-Filtering}: Filtere nach Metadaten (Datum, Kategorie), bevor die Vektorsuche läuft -- reduziert die Suchmenge drastisch.
    \item \textbf{Batch-Indexing}: Embedding-Calls sind teuer. Indexiere in Batches von 100+ Chunks statt einzeln.
\end{itemize}

\begin{lstlisting}[language=Python, caption=Batch-Indexing für bessere Performance]
def batch_index(documents: list[str],
                embedder, vector_store,
                batch_size: int = 100):

    for i in range(0, len(documents), batch_size):
        batch = documents[i:i + batch_size]
        embeddings = embedder.embed(batch)
        vector_store.add_embeddings(batch, embeddings)
        print(f"Indexed {i + len(batch)} / {len(documents)} chunks")
\end{lstlisting}

% ===========================================================%
% SICHERHEIT
% ===========================================================%
\section{Sicherheit}

\subsection{RAG-spezifische Sicherheitsrisiken}

\begin{itemize}[leftmargin=*]
    \item \textbf{Indirect Prompt Injection}: Ein Dokument im Vector Store kann schädliche Anweisungen enthalten. Wenn es retrieved wird, gelangen sie in den LLM-Kontext.
    \item \textbf{Dokumenten-Leakage}: Durch geschickte Queries können Angreifer Dokumente aus dem Vector Store extrahieren, die sie nicht sehen sollten.
    \item \textbf{Tenant-Isolation}: In Multi-Tenant-Systemen müssen die Chunks pro Mandant getrennt sein -- sonst sieht Tenant A die Daten von Tenant B.
    \item \textbf{Poisoning}: Ein Angreifer lädt schädliche Dokumente hoch, die den Index manipulieren.
\end{itemize}

\subsection{Gegenmaßnahmen}

\begin{lstlisting}[language=Python, caption=Tenant-Isolation in RAG]
class TenantAwareRetriever:
    def __init__(self, vector_store, tenant_id: str):
        self.store = vector_store
        self.tenant_id = tenant_id

    def search(self, query: str, top_k: int = 5) -> list[dict]:
        # Metadaten-Filter stellt sicher: nur Chunks dieses Tenants
        return self.store.similarity_search(
            query,
            k=top_k,
            filter={"tenant_id": self.tenant_id},
        )

    def add_document(self, content: str, embedding: list[float]):
        self.store.insert(
            content=content,
            embedding=embedding,
            metadata={"tenant_id": self.tenant_id},
        )
\end{lstlisting}

\textbf{Weitere Sicherheitsmaßnahmen:}
\begin{itemize}[leftmargin=*]
    \item \textbf{Input-Sanitization}: Dokumente vor dem Indexieren auf schädliche Anweisungen scannen
    \item \textbf{Output-Filter}: LLM-Response auf Leaked Content prüfen
    \item \textbf{Access Control}: Wer darf welche Dokumente indexieren? Wer darf welche Collections durchsuchen?
    \item \textbf{Audit-Log}: Welcher User hat was retrieved? (Wichtig für DSGVO-Compliance)
\end{itemize}

% ===========================================================%
% ZUSAMMENFASSUNG
% ===========================================================%
\section{Zusammenfassung}

\begin{itemize}[leftmargin=*]
    \item RAG löst das Wissens- und Aktualitätsproblem von LLMs, indem es externe Dokumente als Kontext in den Prompt einfügt
    \item Die vier Phasen: Chunking $\rightarrow$ Embedding $\rightarrow$ Vector Storage $\rightarrow$ Retrieval + Generation
    \item Die Chunking-Strategie ist der wichtigste Hebel für Retrieval-Qualität -- Recursive Chunking (512 Tokens, 64 Overlap) ist der empfohlene Startpunkt
    \item Voyage, Cohere und OpenAI text-embedding-3-small sind die führenden Embedding-Modelle für Produktion
    \item pgvector (PostgreSQL) und ChromaDB sind die einfachsten Einstiegspunkte; Pinecone und Qdrant für Scale
    \item Erweiterte Techniken (HyDE, Multi-Query, RAPTOR, CRAG, Self-RAG) lösen spezifische Schwächen der naiven RAG
    \item Hit Rate, MRR, NDCG und Faithfulness sind die zentralen Metriken für RAG-Qualität
    \item Security: Indirect Prompt Injection, Tenant-Isolation und Access Control sind produktionskritisch
\end{itemize}

% ===========================================================%
% MERKE-BOX
% ===========================================================%
\begin{merke}
\begin{itemize}
    \item \textbf{RAG ist kein Fine-Tuning-Ersatz}: RAG liefert Kontext, Fine-Tuning formt den Stil. Beide ergänzen sich.
    \item \textbf{Chunking ist der größte Hebel}: Falsche Chunk-Größe oder fehlender Overlap ruinieren jedes Retrieval.
    \item \textbf{Hybride Suche beat Embedding-Only}: Combine BM25 + Embedding für robuste Ergebnisse.
    \item \textbf{Metriken vor Optimierung}: Miss Hit Rate und Faithfulness, bevor du an Chunking oder Modell schraubst.
    \item \textbf{Every Document is an Injection Vector}: Jedes Dokument im Vector Store kann Prompt-Injection enthalten.
\end{itemize}
\end{merke}

% ===========================================================%
% PRAXISPROJEKT
% ===========================================================%
\section{Praxisprojekt}

\subsection{Aufgabe: RAG-System für eine Produkt-Wissensdatenbank}

Baue ein RAG-System, das Kundenanfragen zu einem fiktiven E-Commerce-Produktkatalog beantwortet.

\textbf{Anforderungen:}
\begin{itemize}[leftmargin=*]
    \item Erstelle einen kleinen Produktkatalog (10 Produkte, je 50--200 Wörter Beschreibung) als JSON
    \item Implementiere eine vollständige RAG-Pipeline: Chunking $\rightarrow$ Embedding $\rightarrow$ pgvector oder ChromaDB $\rightarrow$ Retrieval $\rightarrow$ LLM
    \item Nutze hybride Suche (Keyword + Embedding) für bessere Ergebnisse
    \item Füge einen Cross-Encoder Re-Ranker hinzu
    \item Implementiere HyDE als optionale Verbesserung
    \item Miss Hit Rate @ 5 und Faithfulness auf 20 Testfragen
    \item Füge Tenant-Isolation hinzu (zwei fiktive Mandanten)
\end{itemize}

\textbf{Testdaten:} 10 Produktbeschreibungen + 20 Testfragen (10 einfach, 5 komplex, 5 die kein Produkt betreffen).

\textbf{Erfolgskriterium:} Hit Rate @ 5 > 0,9, Faithfulness > 0,85. Keine Daten von Tenant A landen in Tenant Bs Antworten. HyDE verbessert Hit Rate um mindestens 5\,\%.

% ===========================================================%
% INTERVIEWFRAGEN
% ===========================================================%
\section{Interviewfragen}

\begin{enumerate}[leftmargin=*]
    \item \textbf{Erkläre die vier Phasen von RAG.}

    \textbf{Antwort:} (1) Indexing -- Dokumente chunken und embedden. (2) Retrieval -- ANN-Suche im Vector-Index. (3) Augmentation -- Retrieved Chunks in den Prompt einfügen. (4) Generation -- LLM generiert Antwort mit Kontext.

    \item \textbf{Wann verwendest du RAG, wann Fine-Tuning?}

    \textbf{Antwort:} RAG für Faktenwissen und aktuelle Daten (Dokumente, Datenbanken). Fine-Tuning für Stil, Format und Ton (Markenstimme, Antwort-Struktur). In der Praxis oft kombiniert: Fine-Tuning für das Verhalten, RAG für den Inhalt.

    \item \textbf{Was ist der Unterschied zwischen Cosine Similarity und MMR?}

    \textbf{Antwort:} Cosine Similarity findet die ähnlichsten Chunks, die aber oft redundant sind (alle zum selben Thema). MMR (Maximum Marginal Relevance) bestraft Redundanz -- die Top-K-Ergebnisse decken verschiedene Aspekte ab.

    \item \textbf{Wie gehst du vor, wenn ein RAG-System schlechte Antworten liefert?}

    \textbf{Antwort:} Systematische Analyse: (1) Liegt es am Retrieval (Hit Rate messen)? (2) Liegt es an der Chunk-Qualität (Chunks prüfen)? (3) Liegt es am LLM (Prompt ohne RAG testen)? (4) Schrittweise optimieren: Retrieval $\rightarrow$ Chunking $\rightarrow$ Prompt $\rightarrow$ Modell.

    \item \textbf{Was ist HyDE und wann hilft es?}

    \textbf{Antwort:} HyDE generiert eine hypothetische Antwort, embeddet diese statt der Query, und sucht damit. Hilft, wenn die Query kurz ist (2--5 Wörter) und nicht genug semantischen Kontext für gutes Embedding liefert.

    \item \textbf{Wie implementierst du Tenant-Isolation in RAG?}

    \textbf{Antwort:} Metadaten-Filter auf Tenant-ID. Jeder Chunk bekommt beim Indexieren eine Tenant-ID als Metadatum. Der Retriever filtert vor der Similarity-Suche nach der Tenant-ID des anfragenden Users.

    \item \textbf{Welche Chunking-Strategie empfiehlst du für PDF-Handbücher?}

    \textbf{Antwort:} Recursive Chunking auf Kapitel-Ebene (512 Tokens, 64 Overlap) basierend auf Heading-Hierarchie. Die Überschriftenstruktur wird als Metadatum mitgegeben. Das ermöglicht dem LLM, die Quelle korrekt zu zitieren.

    \item \textbf{Was ist ein Cross-Encoder Re-Ranker und warum ist er wichtig?}

    \textbf{Antwort:} Ein Cross-Encoder bewertet Query-Dokument-Paare gemeinsam (im Gegensatz zu Bi-Encodern, die Query und Dokument getrennt embedden). Er ist genauer, aber zu langsam für die initiale Suche. Deshalb: Bi-Encoder für die Vorauswahl, Cross-Encoder für die finale Sortierung.

    \item \textbf{Wie misst du die Qualität eines RAG-Systems?}

    \textbf{Antwort:} Hit Rate @ K (ist die relevante Info in den Top-K?), MRR (wie weit oben?), Faithfulness (wie viel der Antwort ist durch Quellen belegt?), Answer Relevance (beantwortet die Antwort die Frage?). Ein Test-Set mit 50+ annotierten Query-Dokument-Paaren ist die Basis.

    \item \textbf{Dein RAG-System hat 95\,\% Hit Rate, aber die Antworten sind trotzdem schlecht. Was ist los?}

    \textbf{Antwort:} Das Retrieval ist gut, aber die Augmentation/Generation ist das Problem. Mögliche Ursachen: (1) Der Prompt sagt nicht klar, dass nur der Kontext verwendet werden soll. (2) Das LLM halluziniert trotz Kontext. (3) Die Chunks enthalten zu viele irrelevante Details. (4) Der Re-Ranker sortiert nicht optimal.
\end{enumerate}

% ===========================================================%
% WEITERFÜHRENDE RESSOURCEN
% ===========================================================%
\section{Weiterführende Ressourcen}

\subsection{Dokumentationen}

\begin{itemize}[leftmargin=*]
    \item LangChain RAG Guide: \href{https://python.langchain.com/docs/tutorials/rag/}{python.langchain.com/docs/tutorials/rag/}
    \item LlamaIndex RAG: \href{https://docs.llamaindex.ai/en/stable/}{docs.llamaindex.ai}
    \item Pinecone Docs: \href{https://docs.pinecone.io}{docs.pinecone.io}
    \item pgvector GitHub: \href{https://github.com/pgvector/pgvector}{github.com/pgvector/pgvector}
\end{itemize}

\subsection{Tools}

\begin{itemize}[leftmargin=*]
    \item \textbf{LangChain}: \href{https://python.langchain.com}{python.langchain.com} -- Framework für RAG-Pipelines
    \item \textbf{LlamaIndex}: \href{https://www.llamaindex.ai}{llamaindex.ai} -- Spezialisiert auf Indexing und Retrieval
    \item \textbf{RAGAS}: \href{https://docs.ragas.io}{docs.ragas.io} -- RAG-Evaluierungs-Framework
    \item \textbf{Unstructured.io}: \href{https://unstructured.io}{unstructured.io} -- Dokumenten-Parsing für RAG
    \item \textbf{Voyage AI}: \href{https://voyageai.com}{voyageai.com} -- Domain-spezifische Embeddings
\end{itemize}

\subsection{Literatur}

\begin{itemize}[leftmargin=*]
    \item Lewis et al. (2020): "Retrieval-Augmented Generation for Knowledge-Intensive NLP Tasks" -- grundlegendes RAG-Paper
    \item Gao et al. (2023): "Retrieval-Augmented Generation for Large Language Models: A Survey" -- umfassender RAG-Survey
    \item Gao et al. (2022): "Precise Zero-Shot Dense Retrieval without Relevance Labels" -- HyDE-Paper
    \item Sarti et al. (2024): "RAGAS: Automated Evaluation of Retrieval Augmented Generation"
    \item Microsoft (2024): "GraphRAG: Unlocking LLM Discovery on Narrative Private Data"
\end{itemize}

\textbf{Nächster Schritt:} Mit RAG-Wissen ausgestattet geht es weiter zu AI Agents -- der nächsten Evolutionsstufe der LLM-Integration.
